\uuid{0auy}
\titre{Limite en l'infini}
\auteur{bodin ; hanani ; mzari}
\organisation{exo7}
\licence{CC BY-NC-SA 4.0}
\datecreate{2020-01-20}
\niveau{L1}
\module{Analyse}
\chapitre{Continuité, limite et étude de fonctions réelles}
\sousChapitre{Limite de fonctions}
\difficulte{3}
\theme{Limites des fonctions réelles, limite en l'infini}
\exo7qcmid{lille-1.limites.027}

\contenu{
\begin{question}
\qclassification{123.03}
\qtags{motcle=limite en l'infini}
Soit  $f(x)=x+\sqrt[5]{1-x^5}$. Quelles sont les assertions vraies ?
\begin{answers}

    \good{$\lim_{x\to +\infty} f(x)=0$}

    \bad{$\lim_{x\to +\infty } f(x)=+\infty$}

    \bad{$\lim_{x\to -\infty } f(x)=-\infty$}

    \good{$\lim_{x\to -\infty} f(x)=0$}
\end{answers}
\begin{explanations}
En utilisant l'égalité : $a^5+b^5=(a+b)(a^4-a^3b+a^2b^2-ab^3+b^4)$, on pourra multiplier $f$ par l'expression conjuguée de $x+\sqrt[5]{1-x^5}$. On obtient :
 $f(x)=\big[x^4-x^3\sqrt[5]{1-x^5} +x^2(\sqrt[5]{1-x^5})^2-x(\sqrt[5]{1-x^5})^3+(\sqrt[5]{1-x^5})^4\big]^{-1}$.
\end{explanations}
\end{question}
}
